一个质量为 $m>0$ 的非相对论粒子被限制在一维无限深方势阱 $0<x<L$ 内，$L>0$；左右端边界条件 $\psi(0)=\psi(L)=0$。右端是一块可非常缓慢移动的理想量子活塞，忽略其有限运动速度引起的非绝热跃迁。取 $\hbar,k_B>0$，势阱内没有其他势。
(1) 从定态薛定谔方程求所有归一化能级波函数与能量；(2) 定义第 $n$ 定态对右侧活塞的广义压力（即一维力） $F_n=-\partial E_n/\partial L$，求其表达式并解释正负号；(3) 粒子起初处于第 $n$ 定态，右侧活塞绝热缓慢由 $L_i$ 移到 $L_f>L_i$，求输出功，验证量子数保持不变与能量守恒；(4) 若粒子在固定温度 $T$ 下与热库平衡且 $k_BT\gg\hbar^2/(mL^2)$，利用正则配分函数的连续近似求平均活塞力与平均能量，恢复一维经典理想气体的极限。
